% Generated draft: requires mathematical and source review.
\documentclass[11pt]{article}
\usepackage[margin=25mm]{geometry}
\usepackage{fontspec}
\setmainfont{Latin Modern Roman}
\usepackage{amsmath,amssymb,mathtools,unicode-math}
\setmathfont{Latin Modern Math}
\usepackage{enumitem,needspace,xurl,hyperref}
\setlength{\parindent}{0pt}
\setlength{\parskip}{5pt}
\setlength{\emergencystretch}{3em}
\newcommand{\N}{\mathbb{N}}
\newcommand{\Z}{\mathbb{Z}}
\newcommand{\Q}{\mathbb{Q}}
\newcommand{\R}{\mathbb{R}}
\newcommand{\C}{\mathbb{C}}
\newcommand{\F}{\mathbb{F}}
\newcommand{\A}{\mathbb{A}}
\newcommand{\PP}{\mathbb{P}}
\newcommand{\E}{\mathbb{E}}
\newcommand{\eps}{\varepsilon}
\newcommand{\dd}{\,\mathrm{d}}
\begin{document}
\section*{UnsolvedMath ID: 30003044}
\textbf{Source problem code:} catalogue entry
\subsection{Definitions and mathematical statement}
The supplied record asks for the following equality of Kato-cohomology groups. Its characteristic and notation conventions are not explicitly stated. The formulation below uses the conventional characteristic-two interpretation; these conventions, including the ranges of the integer parameters, require confirmation against the cited report.

Let \(F\) be a field of characteristic \(2\), let \(m,n\in\Z_{\geq 0}\), let \(a_1,\ldots,a_m\in F^\times\), and let \(b\in F\). Write \(\langle 1,a_i\rangle\) for the symmetric bilinear form with diagonal matrix \(\operatorname{diag}(1,a_i)\), and write \([1,b]\) for the quadratic form \((x,y)\mapsto x^2+xy+by^2\). Suppose that the quadratic Pfister form
\[
Q=\langle 1,a_1\rangle\otimes\cdots\otimes\langle 1,a_m\rangle\otimes[1,b]
\]
is anisotropic, meaning that \(Q(v)=0\) implies \(v=0\). Let \(F(Q)\) denote the function field of the projective quadric defined by \(Q=0\).

For any field \(K\) of characteristic \(2\), let \(\Omega_K^r=\bigwedge_K^r\Omega^1_{K/\F_2}\) be the space of absolute differential \(r\)-forms, with \(\Omega_K^0=K\) and \(\Omega_K^{-1}=0\). Write \(\mathrm{dlog}(u)=\dd u/u\) for \(u\in K^\times\). The standard Kato operator
\[
\wp_r:\Omega_K^r\longrightarrow\Omega_K^r/\dd\Omega_K^{r-1}
\]
is the additive map characterized on logarithmic expressions by
\[
\wp_r\!\left(x\,\mathrm{dlog}(u_1)\wedge\cdots\wedge\mathrm{dlog}(u_r)\right)
=(x^2-x)\,\mathrm{dlog}(u_1)\wedge\cdots\wedge\mathrm{dlog}(u_r)
\pmod{\dd\Omega_K^{r-1}}.
\]
Define
\[
\nu_K(r)=\ker(\wp_r),\qquad
H_2^{r+1}(K)=\operatorname{coker}(\wp_r).
\]
In degree zero, the logarithmic product is empty, so \(\wp_0(x)=x^2-x\) and \(\nu_K(0)=\F_2\). For an extension \(L/K\), define
\[
H_2^{r+1}(L/K)=\ker\!\left(H_2^{r+1}(K)\longrightarrow H_2^{r+1}(L)\right),
\]
where the map is induced by extension of scalars on differential forms.

Under these conventions, is it true for every such choice of data that
\[
H_2^{n+1}(F(Q)/F)=
\begin{cases}
\left\{\left[b\,\mathrm{dlog}(a_1)\wedge\cdots\wedge\mathrm{dlog}(a_m)\wedge\eta\right]:\eta\in\nu_F(n-m)\right\},&n\geq m,\\[4pt]
\{0\},&n<m,
\end{cases}
\]
with brackets denoting classes in \(H_2^{n+1}(F)\)? Empty products are interpreted as \(1\) if \(m=0\); whether that boundary case is included in the report's intended question is unresolved from the supplied record.

The corrected statement supplied in the registry uses \(H_2^{n+1}\). The field named \texttt{original\_statement} instead prints \(H^{2n+1}\) and leaves the symbol unspecified. This draft follows the supplied corrected version, whose verification is reported by the registry and has not been independently checked here.
\subsection{Short English statement}
For an anisotropic quadratic Pfister form, ask whether the kernel of Kato cohomology after passage to its quadric's function field consists exactly of the displayed differential-symbol classes when \(n\geq m\), and vanishes when \(n<m\). This draft uses the conventional characteristic-two definitions; the source record omits explicit characteristic and parameter-range hypotheses. The supplied status is open, subject to its dated curation and human mathematical review.
\subsection{Sources}
\begin{itemize}
\item Local UnsolvedMath registry, permanent ID \(30003044\), title \emph{Kato Cohomology Kernels of Quadratic Pfister Forms}, source problem code \texttt{catalogue entry}, in the Oberwolfach Reports Open Problems collection. The record gives proposed year \(2016\) and names no individual proposer.
\item Supplied report citation: \emph{Algebraic Cobordism and Projective Homogeneous Varieties} (2016), Oberwolfach Reports, pp.~215--265, DOI \texttt{10.4171/owr/2016/5}. Original OWR report: \url{https://doi.org/10.4171/owr/2016/5}. The supplied background describes the reports as records of mathematical workshops and their open-problem discussions.
\item Supplied related literature: Aravire--Laghribi--O'Ryan kernel computations, \url{https://doi.org/10.1016/j.jalgebra.2015.11.022}. No further bibliographic details are supplied.
\item Literature assessment dated exactly \texttt{2026-08-21}: supplied dated curation, not independently verified. It labels the problem open and states that the report records positive special cases and a conditional theorem, while the displayed equality for general Pfister data remains a question; the related kernel computations do not prove the full statement. The record explicitly requests independent verification of time-sensitive status and cited claims.
\item The registry labels its statement \texttt{corrected\_verified} and reports that the official passage was checked and rewritten as a self-contained problem preserving its mathematical content. This is supplied provenance, not an independent verification by this draft.
\end{itemize}
\end{document}
